\VerdierPhysicalPageAnchor{0069}%
\VerdierSemanticAnchor{ch1-2-2-2-continuation}%
forment un complexe $m$-uple de $C$.

\VerdierSemanticAnchor{ch1-2-2-2-morphism}
Soit maintenant $f\colon Y\longrightarrow Y'$ un morphisme de complexes
$n$-uples de $C$, $\varphi$-sommables. Posons :

\VerdierSemanticAnchor{ch1-2-2-2-eq6}
\begin{equation}
  \left(\int_\varphi f\right)^\tau
  =\bigoplus_{\varphi(\sigma)=\tau}f^\sigma\colon
  \left(\int_\varphi Y\right)^\tau
  \longrightarrow
  \left(\int_\varphi Y'\right)^\tau.
  \tag{2.2.2.6}
\end{equation}

\VerdierSemanticAnchor{ch1-2-2-2-functoriality}
Il est clair que les $\left(\int_\varphi f\right)^\tau$ commutent aux
différentielles de $\int_\varphi Y$ et de $\int_\varphi Y'$ et que, par
suite, une telle famille définit un morphisme de complexes $m$-uples
$\int_\varphi f\colon\int_\varphi Y\longrightarrow\int_\varphi Y'$.
Il est clair que
$\int_\varphi fg=\int_\varphi f\,\int_\varphi g$ et que
$\int_\varphi\operatorname{id}_Y
=\operatorname{id}_{\int_\varphi Y}$. Par suite $\int_\varphi$ est un
foncteur de la catégorie des complexes $n$-uples $\varphi$-sommables dans
la catégorie des complexes $m$-uples. Ce foncteur est additif.

\VerdierSemanticAnchor{ch1-2-2-3}
\noindent\textbf{2.2.3.}\quad
Soient $\varphi\colon[n]\longrightarrow[m]$ une application, $\tau$ un
élément de $\mathbb Z[m]$ et $i$ un élément de $[m]$. Soient $E_1$
(resp. $E_2$) les ensembles :

\VerdierSemanticAnchor{ch1-2-2-3-fibers}
\[
  \begin{aligned}
    E_1&=\{\sigma_1\in\mathbb Z[n]\mid\varphi(\sigma_1)=\tau\},\\
    \text{(resp. }E_2
      &=\{\sigma_2\in\mathbb Z[n]\mid
          \varphi(\sigma_2)=\tau+e_i\}\text{)}.
  \end{aligned}
\]
Soit enfin $Y$ un complexe $n$-uple de $C$. Posons, pour
$\sigma_1\in E_1$, $\sigma_2\in E_2$ :

\VerdierSemanticAnchor{ch1-2-2-3-eq1}
\begin{equation}
  u_{\sigma_2,\sigma_1}(Y)
  =\begin{cases}
    \varepsilon(\tau,i)\varepsilon(\sigma_1,j)d_Y^{\sigma_1,j},
      &\text{si }\sigma_2-\sigma_1=e_j,\\
    0,&\text{sinon},
  \end{cases}
  \tag{2.2.3.1}
\end{equation}
d'où une matrice de type $E_2\times E_1$ :

\VerdierSemanticAnchor{ch1-2-2-3-eq2}
\begin{equation}
  M(\tau,i,Y)
  =\bigl(u_{\sigma_2,\sigma_1}(Y)\bigr)_
    {\sigma_2\in E_2,\,\sigma_1\in E_1}.
  \tag{2.2.3.2}
\end{equation}

\VerdierSemanticAnchor{ch1-2-2-3-matrix-identification}
Lorsque $Y$ est $\varphi$-sommable, cette matrice est la matrice du
morphisme $d_{\int_\varphi Y}^{\tau,i}$.

\VerdierSemanticAnchor{ch1-2-2-3-extensions}
Remarquons que dans le cas général, cette matrice ne possède qu'un nombre
fini de termes non nuls sur chaque ligne et sur chaque colonne. Il y a donc
\emph{a priori} au moins deux extensions possibles du foncteur
$\int_\varphi$ aux complexes quelconques lorsque, dans $C$, les produits
dénombrables et les sommes directes
